\documentclass{ou}
\begin{document}
\thispagestyle{empty}
\begin{center}
\shortstack{\FGAVO}
\vfill
\setstretch{2}
\textbf{РЕФЕРАТ}\\
Информационные системы и технологии\\
\setstretch{1}
Технологии создания и обработки графических данных
\vfill
\end{center}
\thispagestyle{empty}
\vspace{-12pt}
\begin{spacing}{1}
\noindent\instructor \hfill И. Р. Нигматуллин\\
\setstretch{2}
\students \hfill ГФ25-01Б\\
\end{spacing}
\begin{center}
    \city
\end{center}
\newpage
\tableofcontents
\newpage
\section{Отличие векторной и растровой графики и области их применения}
\subsection{}
\newpage
\section{Структура и принципы работы цветовых моделей}
Цветовая модель -- это математический (или компьютерный) способ описания цветов. Она не зависит от физических устройств.

Важно различать цветовую модель от пространства, т.к. цветовая модель -- это соответствие между некоторыми неограниченными действительными числами и цветами, а цветовое пространство -- это набор ограничений для этих чисел. Многие цветовые модели способны воспроизводить бесконечную палитру цветов, но соответствующие цветовые пространства часто ограничивают параметры определённым диапазоном (например, в модели RGB насыщенный красный цвет может быть представлен как (150 \verb|%, 0 %, 0 %)|, но конкретное цветовое пространство, связанное с моделью, может ограничивать значения красного в диапазоне от 0 \verb|% до 100 %| и, таким образом, не сможет отобразить этот цвет)

Чтобы объяснить шестилетнему ребёнку, что такое цветовое пространство, я мог бы сказать что-то вроде этого. У вас есть 44 цветных карандаша, и вы можете использовать только эти 44 цветных карандаша для создания изображения.

Сущестует два типа смешения цветов: субтрактивный и аддитивный

Аддитивное смешение работает со светом. Основные цвета -- красный, зеленый и синий (RGB). При их смешении лучи света складываются. Совмещение всех трех цветов в полной интенсивности дает белый свет. RGB модель используют все излучающие свет устройства: мониторы, телевизоры и телефоны.

Субтрактивное смешение работает с отраженным светом, то есть с красками, чернилами или пигментами. Основные цвета -- голубой, пурпурный и желтый (CMY). При их смешении каждый пигмент поглощает (вычитает) определенную часть спектра. Теоретически смешение трех основных красок должно давать черный цвет, но на практике получается грязно-коричневый, поэтому в модель CMYK добавлен отдельный черный цвет (K). Эта модель является стандартом для цветной печати.
\subsection{RGB}
Каждый пиксель на экране состоит из трех крошечных светодиодов или ячеек, излучающих свет определенной интенсивности. Интенсивность каждого цвета измеряется в диапазоне от 0 до 255, где 0 означает полное отсутствие цвета, а 255 - его максимальную яркость. 

Например, чистый красный цвет записывается как RGB(255,0,0), а белый - RGB(255,255,255). При наложении всех трех базовых цветов максимальной интенсивности получается белый цвет, а при их отсутствии - черный. Такой принцип работы повторяет механизм восприятия цвета человеческим глазом. 

Для наглядности его можно представить в виде куба

При работе с цветовой моделью RGB каждый цвет можно закодировать двумя способами - через десятичные или шестнадцатеричные значения. В десятичной системе интенсивность каждого базового цвета задается числом от 0 до 255, где 0 означает отсутствие цвета, а 255 - его максимальную яркость.

Шестнадцатеричная запись использует символы от 0 до 9 и буквы от A до F. Каждый базовый цвет кодируется двумя символами, что дает 256 градаций (от 00 до FF). Запись начинается с символа \verb|#|, за которым следуют 6 знаков - по два на красный, зеленый и синий. Например, \verb|#|FF0000 задает чистый красный цвет. 
\subsection{CMYK}
Цветовая модель CMYK включает четыре базовых компонента, каждый из которых выполняет специфическую функцию при создании полноцветных изображений. Cyan (голубой) отвечает за холодные оттенки и формирует основу для создания зеленых и синих тонов. Magenta (пурпурный) позволяет получать красные и розовые оттенки, участвуя в создании теплой гаммы изображения. Yellow (желтый) служит базой для воспроизведения солнечных и теплых тонов.

Black (черный) компонент в схеме CMYK выполняет две основные задачи: усиливает контрастность и глубину изображения, а также обеспечивает четкость текста и тонких линий. При печати черный краситель наносится последним, что позволяет добиться максимальной насыщенности темных участков.

Процентное соотношение компонентов определяет итоговый цвет. Например, для получения травянисто-зеленого используется комбинация 100\verb|%| cyan и 100\verb|%| yellow. Насыщенный фиолетовый создается при смешении 100\verb|%| magenta и 50\verb|%| cyan. Точность передачи цветов зависит от качества красителей и калибровки печатного оборудования. 
\subsection{HSB}
HSB представляет собой цветовую модель, разработанную для интуитивного описания цветов в терминах, близких к человеческому восприятию. В отличие от технических моделей вроде RGB или СМУК, HSB напрямую оперирует понятиями, которые мы используем в повседневной речи, когда говорим о цветах. Три фундаментальных параметра HSB формируют полные координаты цвета в этой системе: 

- Цветовой тон (Hue);

- Насыщенность (Saturation);

- Яркость (Brightness).

Цветовой тон -- это то, что мы обычно называем цветом в повседневной речи. Насыщенность -- показывает интенсивность или чистоту цвета. При уменьшении насыщенности цвет становится более приглушённым. Яркость определяет, насколько светлым или тёмным является цвет. 

Взаимодействие этих трёх параметров создаёт трёхмерное пространство, в котором можно точно определить любой видимый цвет. Важно понимать, что HSB -- это не независимая система кодирования цветов, а скорее альтернативный способ представления тех же цветов, которые можно описать в RGB или других моделях.
\newpage
\section{Популярные форматы графических файлов и способы сжатия изображений}
\subsection{}
\newpage
\section{Основы редактирования графических данных в прикладных программах}
\subsection{}
\newpage
\section{Цифровая обработка изображений и улучшение их качества}
Цифровая обработка изображений (ЦОИ) относится к цифровой обработке сигналов, где сигналом служит изображение и представляет собой самостоятельную область знаний, которая быстро развивается и охватывает большой спектр методов, которые имеют очень широкое применение. Методы обработки изображений играют значительную роль в научных исследованиях, промышленности, медицине, космических исследованиях и информационных системах. Примерами применения этих методов могут служить цифровая передача изображений с космических кораблей, повышение четкости изображений, создаваемых электронным микроскопом, коррекция искажений изображений, принимаемых из космоса, автоматический анализ характера местности, исследование природных ресурсов по фотоснимкам, передаваемым со спутников Земли, формирование и улучшение качества биологических и медицинских изображений, автоматическое составление карт по аэрофотоснимкам. 

Повышение качества изображений достигается двумя видами обработки изображений: реставрацией изображений и их улучшением. Под реставрацией обычно понимают процедуру восстановления или оценивания элементов изображения, целью которой является коррекция искажений и наилучшая аппроксимация идеального неискаженного изображения. Для улучшения изображений используется комплекс операций, призванных улучшить восприятие изображения наблюдателем или же преобразовать его в другое изображение, более удобное для машинной обработки. 

При машинной обработке улучшение изображения тесно связано с задачей извлечения информации. Например, система улучшения изображений производит подчеркивание границ исследуемого изображения путем высокочастотной фильтрации. На основе обработанного изображения выполняется операция прослеживания контура объекта, определяется его форма и размеры. В этом случае система улучшения изображения используется для того, чтобы подчеркнуть важнейшие признаки исходного изображения и облегчить задачу извлечения информации. Манипуляции с данными видимых изображений -- это одно из основных применений обработки изображений. Обработки изображений может быть использована для того, чтобы получить видимое изображение цифровых данных, усиленных определенным образом для выделения некоторых аспектов данных. Такие типы обработки изображения используются, например, в медицинском оборудовании, гидролокаторах, радарах, ультразвуковом оборудовании, теплочувствительном оборудовании и т.д.

Цель обработки изображений:

- Преобразование изображений, ориентированное на визуальное восприятие человеком (в основном повышение качества, коррекция цвета, контраста, исправления мелких помех или искажение изображения);

- Преобразование изображений, ориентированное на их автоматизированный анализ (например, фильтрация, математическая морфология, сегментация на области определенных классов, выделение объектов и т. п.).

Среди задач обработки изображений можно выделить две: улучшение и восстановление изображений.

Главная цель улучшения заключается в таком изменении изображения, чтобы результат оказался наиболее подходящим с точки зрения конкретного применения. Слово конкретное является важным, поскольку, например, метод, будучи полезным для улучшения рентгеновских изображений, как правило, не подходит для обработки цветных фотографий, и наоборот.
\subsection{Применение цифровой обработки изображений}
Можно использовать такие методы, как шумоподавление, устранение размытия и коррекция цвета. Алгоритмы можно использовать для идентификации и извлечения определенных объектов или функций на изображении, таких как лица, транспортные средства или текст. Методы сжатия изображений -- для уменьшения размера файла при сохранении визуального качества. Цифровая обработка изображений широко используется в медицинских приложениях, таких как анализ рентгеновских изображений и изображений МРТ, для диагностики и лечения заболеваний. Цифровая обработка изображений -- в беспилотных автомобилях, системах безопасности и робототехнике, также для наложения виртуальных элементов на изображения реального мира.
\subsection{Преимущества цифровой обработки изображений}
Алгоритмы цифровой обработки изображений могут улучшить визуальное качество изображений, делая их более четкими, резкими и информативными.

Цифровая обработка изображений может автоматизировать многие задачи, связанные с изображениями, такие как распознавание объектов, обнаружение закономерностей и измерение.

Алгоритмы цифровой обработки изображений могут обрабатывать изображения намного быстрее, чем люди, что позволяет анализировать большие объемы данных за короткий промежуток времени.

Алгоритмы цифровой обработки изображений могут обеспечить более точные результаты, чем люди, особенно для задач, требующих точных измерений или количественного анализа.
\subsection{Недостатки цифровой обработки изображений}
Некоторые алгоритмы обработки цифровых изображений требуют больших вычислительных ресурсов и значительных вычислительных ресурсов.

Некоторые алгоритмы цифровой обработки изображений могут давать результаты, которые людям трудно интерпретировать, особенно для сложных или изощренных алгоритмов.

Качество выходных данных алгоритмов цифровой обработки изображений сильно зависит от качества входных изображений. Входные изображения низкого качества могут привести к плохому качеству вывода.

Алгоритмы цифровой обработки изображений имеют ограничения, такие как сложность распознавания объектов в загроможденных или плохо освещенных сценах или неспособность распознавать объекты со значительными деформациями или окклюзиями.

Производительность многих алгоритмов цифровой обработки изображений зависит от качества обучающих данных, используемых для разработки алгоритмов. Низкое качество обучающих данных может привести к плохой работе алгоритма.
\subsection{Возможности цифровой обработки изображений в Matlab}
На сегодняшний день система Matlab, в частности пакет прикладных программ Image Processing Toolbox, является наиболее мощным инструментом для моделирования и исследования методов обработки изображений. Он включает большое количество встроенных функций, реализующих наиболее распространенные методы обработки изображений.
\subsection{Геометрические преобразования изображений}
К наиболее распространенным функциям геометрических преобразований относится кадрирование изображений (imcrop), изменение размеров (imresize) и поворот изображения (imrotate).

Суть кадрирования состоит в том, что функция imcrop позволяет с помощью мыши в интерактивном режиме вырезать часть изображения и поместить ее в новое окно просмотра.

Функция изменения размеров изображения imresize позволяет, используя специальные методы интерполяции, изменять размер любого типа изображения.

В пакете Image Processing Toolbox существует функция imrotate, которая осуществляет поворот изображения на заданный угол.

Таким образом, приведенные выше функции позволяют поворачивать, вырезать части, масштабировать, т.е. работать с целым массивом изображения.
\subsection{Анализ изображений}
Для работы с отдельными элементами изображений используются такие функции как imhist, impixel, mean2, corr2 и другие.

Одной из наиболее важных характеристик является гистограмма распределения значений интенсивностей пикселей изображения, которую можно построить с помощью функции imhist.

Довольно часто при проведении анализа изображений возникает необходимость определить значения интенсивностей некоторых пикселей. Для этого в интерактивном режиме можно использовать функцию impixel.

Еще одной широко применяемой функцией является функция mean2 -- она вычисляет среднее значение элементов матрицы.
\subsection{Фильтрация изображений}
Пакет Image Processing Toolbox обладает очень мощным инструментарием по фильтрации изображений. Среди множества встроенных функций, которые решают задачи фильтрации изображений, особо следует выделить fspecial, ordfilt2, medfilt2.

Функция fspecial является функцией задания маски предопределенного фильтра. Эта функция позволяет формировать маски:

- Высокочастотного фильтра Лапласа;

- Фильтра, аналогичного последовательному применению фильтров Гаусса и Лапласа;

- Усредняющего низкочастотного фильтра;

- Фильтра, повышающего резкость изображения.
\subsection{Сегментация изображений}
Сегментация изображений -- это процесс разделения цифрового изображения на несколько сегментов (множество пикселей). Цель -- упростить представление изображения, чтобы его было проще анализировать. 

Среди встроенных функций пакета Image Processing Toolbox, которые применяются при решении задач сегментации изображений, следует выделить qtdecomp, edge и roicolor.

Функция qtdecomp выполняет сегментацию изображения методом разделения и анализа однородности не перекрывающихся блоков изображения.
\subsection{Морфологические операции над бинарными изображениями}
Морфологические операции над бинарными изображениями -- это методы обработки изображений, которые фокусируются на форме и структуре объектов. Они применяются к двоичным изображениям, где интересующие объекты представлены пикселями переднего плана (обычно белого цвета), а фон -- пикселями фона (обычно чёрного цвета). Основная идея -- исследовать изображение с помощью структурирующего элемента (ядра) и изменять значения пикселей на основе их пространственного расположения и формы элемента.

Система Matlab владеет довольно мощным инструментарием морфологической обработки бинарных изображений. Среди основных операций выделим следующие -- эрозия, наращивание, открытие, закрытие, удаление изолированных пикселей, построение скелета изображения и т.п.
\subsection{Методы улучшения качества изображений}
Основываясь на существующей классификации и атрибутах качества изображения, распределим методы улучшения изображения на пять категорий. Классифицирование выполнено в соответствии с атрибутами качества изображения, которые улучшаются, потому что улучшение обычно выполняется с определенной целью. Для того, чтобы область исследований по улучшению изображения достигла своих конкретных ожиданий улучшения, рассмотрим методы в следующих категориях:

1) Повышение контрастности изображения;

2) Повышение резкости изображения;

3) Цветокоррекция изображения;

4) Де-искусственное изображение;

5) Улучшение изображения для нескольких атрибутов качества.

Гистограмма -- это диаграмма, построенная в столбиковой форме, в которой величина показателя изображается графически в виде столбика. На основе анализа гистограммы можно судить о яркостных искажениях изображения, т.е. сказать о том, является ли изображение затемненным или засветленным. Известно, что в идеале на цифровом изображении в равном количестве должны присутствовать пиксели со всеми значениями яркостей, т.е. гистограмма должна быть равномерной. Перераспределение яркостей пикселей на изображении с целью получения равномерной гистограммы выполняет метод эквализации, который в системе Matlab реализован в виде функции histeq. Улучшение изображений путем выравнивания гистограммы -- это процесс, в котором пытаются достичь равномерности распределения яркостей обработанного изображения. Процедура выравнивания гистограммы состоит из следующих действий:

- Вычисляется гистограмма распределения яркостей элементов изображения H(L);

- Строится нормированная кумулятивная гистограмма CH(L);

- Формируется новое изображение L' = R*CH(L).

Это преобразование эффективно для улучшения визуального качества низко контрастных деталей. Описанный метод преобразования гистограммы может быть глобальными, то есть использовать информацию обо всем изображении, и скользящими, когда для преобразования используются локальные области изображения. Глобальные методы гистограммных преобразований являются, в сущности, табличными методами, основное их преимущество состоит в быстродействии. Для более детальной обработки изображений используют скользящие методы. Они обеспечивают качественное контрастирование мелких деталей изображения. Вместе с тем, они более объемны по вычислительным затратам. Поэтому при использовании методов класса гистограммных преобразований нужно искать компромисс между качеством и быстродействием обработки изображений.

Алгоритм построения гистограммы:

- Строим массив, заполняем нулями;

- Цикл, для каждого пикселя;

-	Увеличиваем значение элемента массив[значение] на 1;

-	Конец цикла.

Полученный массив и представляет собой гистограмму, элементы массива -- означают высоты столбиков.

Для цветных изображений процедура построения гистограммы выполняется раздельно для каждого цветового канала, что позволяет проводить независимый анализ цветовых компонент.

Линейные преобразования реализуются посредством аффинных преобразований яркостных характеристик. Базовое преобразование имеет вид L' = k×L + b, где параметр k определяет степень контрастности, а коэффициент b отвечает за коррекцию яркости.

Нелинейные методы коррекции включают:

- Гамма-коррекцию, обеспечивающую нелинейное преобразование динамического диапазона;

- Логарифмические преобразования, эффективно сжимающие широкий динамический диапазон;

- Степенные преобразования, обеспечивающие гибкое управление яркостными характеристиками.

Метод «серый мир» основан на гипотезе о среднесером цвете усредненного изображения. Алгоритм включает:

- Вычисление средних значений по каждому цветовому каналу;

- Определение глобального среднего по всем каналам;

- Коррекцию цветовых компонент с использованием полученных средних значений.

Процедура бинаризации представляет собой этап многих задач компьютерного зрения. Пороговая обработка позволяет разделить пикселы на два класса - объекты и фон.

Адаптивная бинаризация решает проблему неравномерного освещения через:

- Использование скользящего оконного анализа;

- Вычисление локальных порогов для каждого пикселя;

- Применение пространственно-вариативных критериев классификации.

Классификация шумов по их статистическим характеристикам:

- Импульсный шум -- дискретные выбросы яркости;

- Гауссов шум -- непрерывные случайные флуктуации;

- Спектральный шум -- частотно-избирательные искажения.

Методы фильтрации:

- Гауссово сглаживание -- эффективное подавление гауссова шума;

- Медианная фильтрация -- устранение импульсных помех;

- Адаптивные фильтры -- автоматическая подстройка параметров.

Фильтр Unsharp Mask (нерезкая маска) используется для тонкой цветовой коррекции, он является очень полезным инструментом и предназначен для увеличения резкости.

Технология повышения визуальной резкости через:

- Выделение высокочастотных компонент изображения;

- Усиление контраста на границах и контурах;

- Компенсацию оптических искажений.

Побочные эффекты:

- Появление ореолов;

- Усиление шумовых компонент;

- Ступенчатость на плавных градиентах.

Интерполяционные алгоритмы улучшения разрешения:

- Nearest-neighbor -- максимальное быстродействие;

- Bilinear -- сглаживание ступенчатости;

- Bicubic -- сохранение резкости контуров;

- Spline-интерполяция -- минимальные искажения.

Спектральные методы используют:

- Частотный анализ изображений;

- Восстановление высокочастотных компонент;

- Компенсацию аппаратной функции.

Современные подходы включают:

- Машинное обучение для адаптивной обработки;

- Нейросетевые алгоритмы enhancement-а;

- Контекстно-зависимую коррекцию;

- Многомасштабный анализ изображений.

Основные шаги реализации методов адаптивного преобразования локальных контрастов следующие:

1) Для каждого элемента изображения L(i,j) вычисляют значение локального контраста C(i,j) в текущей окрестности W с центром в элементе с координатами (i,j);

2) Вычисляют локальную статистику для текущей скользящей окрестности W;

3) Преобразуют (усиливают) локальный контраст C(i,j), употребляя для этого нелинейные функции и учитывая локальную статистику текущей скользящей окрестности W;

4) Восстанавливают значение яркости изображения L'(i,j) с усиленным локальным контрастом.
\newpage
\section{Организация и хранение графических коллекций в цифровой среде}
Графические коллекции играют важнейшую роль в современном информационном обществе. Они являются ценным источником информации, инструментом коммуникации и важной частью культурного наследия. Эффективная организация и хранение графических коллекций в цифровой среде -- это задача, приобретающая все большую актуальность в связи с ростом объемов цифрового контента и необходимостью обеспечения его доступности, сохранности и использования.
\subsection{Определение и типы графических коллекций}
Цифровая графическая коллекция -- это организованная совокупность графических файлов, представленных в цифровом формате. Это могут быть фотографии, иллюстрации, чертежи, диаграммы, логотипы, иконки и другие виды визуального контента.

Классификация графических коллекций:

- По тематике;

- По источнику;

- По назначению;

- По степени обработки.

Говоря о тематике, сюда относят фотографии природы, исторические фотографии, архитектурные чертежи, медицинские изображения, маркетинговые материалы. Например, коллекция фотографий русской природы, включающая изображения лесов, рек, животных и пейзажей, или коллекция старинных карт европейских городов.

Если рассматривать источник, то коллекции могут состоять из сканированных изображений, цифровых фотографий, векторной графики, сгенерированных компьютером изображений. Например, коллекция оцифрованных слайдов с семейного архива или коллекция логотипов, разработанных дизайнерским агентством.

В разрезе назначения выделяют архивные подборки, коллекции для публикаций, исследования или обучения. Например, архивная коллекция фотографий времен Второй мировой войны или коллекция изображений, используемых для создания рекламных кампаний.

По степени обработки различают необработанные фотографии, отредактированные изображения и готовые к публикации файлы. Например, коллекция RAW-файлов с цифровой камеры или коллекция финальных версий маркетинговых баннеров.
\subsection{Основные принципы организации графических коллекций}
Централизованное и децентрализованное хранение:

- Централизованное хранение;

- Децентрализованное хранение.

При централизованном хранении все графические файлы хранятся в одном месте, обычно на сервере или в облачном хранилище. Это обеспечивает единую точку доступа, упрощает управление правами доступа и обеспечивает более эффективное использование ресурсов.

При децентрализованном хранении графические файлы распределены по нескольким компьютерам или устройствам. Этот подход может быть более гибким и устойчивым к сбоям, но усложняет управление и требует дополнительных мер для обеспечения синхронизации и согласованности.

Метаданные -- это данные о данных. В контексте графических коллекций метаданные представляют собой информацию, описывающую графический файл, такую как:

1) Название (описание изображения);

2) Автор (кто создал изображение);

3) Дата создания (когда было создано изображение);

4) Местоположение (где было сделано изображение);

5) Ключевые слова (теги, описывающие содержание изображения);

6) Описание (более подробное описание изображения);

7) Права (информация о правах на использование изображения).

Использование метаданных позволяет:

- Эффективно искать и находить нужные изображения;

- Автоматизировать процессы обработки и организации изображений;

- Обеспечить долгосрочную сохранность и доступность изображений;

- Защитить авторские права.

DAM -- это программное обеспечение, предназначенное для управления цифровыми активами, включая графические файлы. DAM-системы обеспечивают централизованное хранение, организацию, поиск, управление правами и распространение графических файлов.
\subsection{Методы индексации и каталогизации графических коллекций}
- Иерархическая классификация;

- Тегирование и ключевые слова;

- Автоматическая индексация на основе анализа изображений;

- Контролируемые словари и тезаурусы.

При описании иерархической классификации речь идёт об организации графических файлов в древовидную структуру папок и подпапок. Например, Фотографии/Природа/Леса/Елки. Это простой и интуитивно понятный способ организации, но он может быть неэффективным для больших коллекций и при необходимости поиска изображений по нескольким критериям.

Когда говорится о тегировании и ключевых словах, имеется в виду присвоение графическим файлам тегов или ключевых слов, описывающих их содержание. Например: Фотография\verb|_|заката.jpg с тегами: закат, море, солнце, пляж, отпуск. Этот метод позволяет гибко классифицировать изображения и находить их по различным критериям.

Что касается автоматической индексации на основе анализа изображений, здесь используется программное обеспечение для автоматического анализа изображений и присвоения им тегов или ключевых слов на основе распознавания объектов, сцен и лиц. Например, программа может автоматически определить, что на фотографии изображен человек, собака и дерево. Этот метод может значительно упростить процесс индексации, но требует точных и надежных алгоритмов распознавания.

И наконец, при использовании контролируемых словарей и тезаурусов применяются заранее определенные списки терминов для тегирования изображений. Это обеспечивает согласованность и предотвращает использование разных терминов для описания одного и того же объекта. Например, использование медицинского тезауруса MeSH для индексации медицинских изображений.
\subsection{Форматы графических файлов и их выбор для хранения}
Растровые форматы:

- JPEG;

- PNG;

- TIFF;

- GIF.

JPEG -- популярный формат для фотографий, обеспечивающий хорошее сжатие с потерями качества. Он подходит для хранения изображений для веб-сайтов и социальных сетей.

PNG -- формат для изображений с прозрачным фоном и без потерь качества. Он подходит для хранения графики, логотипов и иконок.

TIFF -- формат для хранения изображений высокого качества без потерь. Он подходит для архивирования и профессиональной печати.

GIF -- формат для анимации и простых изображений.

Векторные форматы:

- SVG;

- EPS;

- AI;

- CDR.

SVG -- открытый формат для векторной графики, поддерживаемый большинством современных браузеров.

EPS -- формат для векторной графики, используемый в полиграфии.

AI -- собственный формат Adobe Illustrator.

CDR -- собственный формат CorelDRAW.

Факторы, влияющие на выбор формата:

- Качество;

- Размер файла;

- Совместимость;

- Сохранение прозрачности.

При выборе качества необходимо обратить внимание на формат, который обеспечивает требуемое качество изображения.

Когда речь идёт о размере файла, важно выбрать формат, обеспечивающий приемлемый объём хранения.

Говоря о совместимости, нужно учитывать форматы, поддерживаемые всеми необходимыми программами и устройствами.

Если речь идёт о сохранении прозрачности, необходимо выбирать формат, который поддерживает прозрачность, например PNG или GIF.
\subsection{Системы хранения графических данных}
- Локальное хранилище;

- Сетевое хранилище (NAS, SAN);

- Облачные хранилища (Amazon S3, Google Cloud Storage, Azure Blob Storage).

Локальное хранилище -- хранение графических файлов на жестком диске компьютера. Это простой и доступный способ хранения, но он не обеспечивает достаточной безопасности и доступности данных.

Сетевое хранилище (NAS, SAN) -- хранение графических файлов на сетевом хранилище, обеспечивающем доступ к данным с нескольких компьютеров. Такой способ хранения более надежный и удобный, но требует дополнительных затрат на оборудование и обслуживание.

Облачное хранилище (Amazon S3, Google Cloud Storage, Azure Blob Storage) -- хранение графических файлов в облаке, которое обеспечивает высокую доступность, надежность и масштабируемость. Этот способ удобный и экономичный, но требует подключения к Интернету и может быть связан с вопросами безопасности.
\subsection{Методы обеспечения целостности и безопасности графических данных}
- Резервное копирование (Backup);

- Системы контроля версий;

- Права доступа и управление пользователями;

- Защита от несанкционированного доступа.

Резервное копирование (Backup) -- создание копий графических файлов и их хранение в другом месте. Это важный метод защиты от потери данных в случае сбоя оборудования или других аварийных ситуаций. Существуют различные стратегии резервного копирования, такие как полное резервное копирование, инкрементное резервное копирование и дифференциальное резервное копирование.

Системы контроля версий -- использование программного обеспечения для отслеживания изменений, внесенных в графические файлы. Это позволяет восстанавливать предыдущие версии файлов и отменять ошибки. Например, Git можно использовать для работы с векторной графикой, хотя это может быть нетривиально.

Права доступа и управления пользователями -- ограничение доступа к графическим файлам для авторизованных пользователей. Существуют различные модели управления правами доступа, такие как ролевая модель и модель доступа на основе атрибутов.

Защита от несанкционированного доступа -- использование паролей, шифрования и других мер для защиты графических файлов. Важно использовать надежные пароли и регулярно их менять.
\subsection{Практические примеры организации графических коллекций в различных сферах}
- Фотоархивы музеев и библиотек;

- Коллекции изображений в дизайнерских агентствах;

- Управление графическими активами в крупных корпорациях.

Говоря о фотоархивах музеев и библиотек, имеют в виду использование DAM-систем для управления большими коллекциями исторических фотографий, карт и других графических материалов. Это обеспечивает доступ к этим материалам для исследователей и широкой публики. Например, оцифрованные коллекции Российской Государственной Библиотеки.

Когда речь идёт о коллекциях изображений в дизайнерских агентствах, подразумевается использование DAM-систем для управления логотипами, иллюстрациями, фотографиями и другими графическими активами, используемыми в проектах дизайна. Это обеспечивает доступ к этим активам для дизайнеров и других сотрудников агентства.

В случае управления графическими активами в крупных корпорациях речь идёт о применении DAM-систем для управления логотипами, фирменным стилем, маркетинговыми материалами и другими графическими активами, используемыми в корпоративной коммуникации. Это позволяет обеспечить согласованность и эффективность использования этих активов.
\subsection{Перспективы развития систем организации и хранения графических коллекций}
Использование искусственного интеллекта и машинного обучения включает автоматическое тегирование, распознавание объектов, лиц и сцен на изображениях, а также использование AI для поиска похожих изображений и автоматической организации коллекций.

Расширение возможностей метаданных предполагает использование семантических технологий и онтологий для более точного и информативного описания графических файлов, с поддержкой различных стандартов метаданных, таких как Dublin Core и IPTC.

Интеграция с другими информационными системами осуществляется через подключение DAM-систем к системам управления контентом (CMS), системам управления взаимоотношениями с клиентами (CRM) и другим корпоративным системам.

Организация и хранение графических коллекций в цифровой среде -- это сложная и многогранная задача, требующая комплексного подхода. Правильный выбор методов индексации, форматов файлов, систем хранения и инструментов управления позволяет обеспечить доступность, сохранность и эффективное использование графических активов. В будущем развитие технологий искусственного интеллекта и семантических технологий позволит значительно упростить и автоматизировать процессы организации и хранения графических коллекций.
\newpage
\section*{Список использованных источников}
\addcontentsline{toc}{section}{Список использованных источников}
\nocite{*}
\bibliographystyle{plainurl}
\bibliography{report7}
\end{document}